\documentclass[10pt,a4paper]{amsart}
\usepackage[margin=19mm]{geometry}
\usepackage{amsmath,amssymb,mathtools}
\usepackage[scr=rsfs,cal=euler]{mathalfa}
\usepackage{xfrac}
\usepackage[unicode,hidelinks,bookmarks=false]{hyperref}
\newcommand{\ie}{i.e.\ }
\newcommand{\cf}{cf.\ }
\newcommand{\resp}{respectively\ }
\newcommand{\Subsection}{\subsection}
\providecommand{\nobreakdash}{\nobreak\hbox{-}}
\setlength{\emergencystretch}{2em}
\multlinegap0pt
\usepackage{bm}
\usepackage{bbm}
\usepackage{dsfont}
\usepackage{xspace}
\usepackage{cancel}
\usepackage{mathtools}
\usepackage{amsthm}
\makeatletter
\@ifundefined{theorem}{\newtheorem{theorem}{Theorem}}{}
\@ifundefined{lemma}{\newtheorem{lemma}[theorem]{Lemma}}{}
\makeatother
\def\cC{\mathscr{C}}
\newcommand{\abs}[1]{\left\vert#1\right\vert}
\newcommand{\ddbar}{\overline\partial}
\newcommand{\ol}{\overline}
\newcommand{\pr}{\partial}
\title[On the second coefficient in the semi-classical expansion of]{On the second coefficient in the semi-classical expansion of Toeplitz Operators}
\author{Chin-Chia Chang, Hendrik Herrmann, Chin-Yu Hsiao}
\date{Source: arXiv:2412.11697v2 (CC BY 4.0)}
\begin{document}
\maketitle

\noindent\emph{Source and licence.} On the second coefficient in the semi-classical expansion of Toeplitz Operators. Version arXiv:2412.11697v2 (CC BY 4.0), distributed under the Creative Commons Attribution 4.0 International licence, as recorded in the arXiv licence metadata for this exact version and shown on the abstract page https://arxiv.org/abs/2412.11697v2. Full rights notice: licenses/arxiv-2412.11697-CC-BY-4.0.txt.\par\smallskip

\noindent\textit{Editorial scope.} Counted unit: the Lemma whose source label is \texttt{e-gue241128yyd} in the article (source0.tex lines 644--669, source label \texttt{e-gue241128yyd}), delivered together with the article's own complete original proof of it (source0.tex lines 671--718, one \texttt{proof} environment of 48 lines). No mathematics is written, repaired or re-derived: statement and proof are the article's own text line for line. Required context retained uncounted: eq:PiFIO (lines 160--163); e-gue241220yydp (lines 624--626). Every definition, display and label of those lines is retained unchanged. Omitted from this package and not redistributed: 1--159, 164--623, 627--643, 670--670, 719--1328. Nothing inside a retained excerpt is omitted or altered: every retained body is delivered exactly as the article's lines are written, so no omission receipt of any kind accompanies this package. Citations: the retained excerpts carry no citation command at all, so no bibliography entry of the article is transcribed and no reference is redistributed. Reference adaptation: none was needed and none was made; every reference inside the delivered unit resolves inside it, so no unresolved reference or citation is produced. Printed slips kept literal: no printed word, symbol, hypothesis or number of the counted statement or of its proof was altered. Layout and class: the article's own class is \texttt{amsart} with packages including geometry, amsmath, amssymb, amsfonts, amsthm, amsrefs, mathrsfs, epic, amsopn, amscd, graphicx, color, transparent, xcolor; the wrapper is the lane's pinned \texttt{amsart} editorial wrapper at 10pt on A4, which restates the theorem-like environments the retained text needs and adds the article's own macro definitions that the retained text needs (\texttt{\textbackslash abs}, \texttt{\textbackslash cC}, \texttt{\textbackslash ddbar}, \texttt{\textbackslash ol}, \texttt{\textbackslash pr}), with extra wrapper packages (bm, bbm, dsfont, xspace, cancel, mathtools). The delivered numbering is the unit's own. Assets: no retained span contains a figure environment, a graphics-inclusion command, a PDF-page inclusion, a table or a TikZ picture; no image file is copied into this package, the package declares no asset dependency and no image file is redistributed with it. Licence: version arXiv:2412.11697v2 of this article is distributed under the Creative Commons Attribution 4.0 International licence, as recorded in the arXiv licence metadata for this exact version and shown on the abstract page https://arxiv.org/abs/2412.11697v2. A full rights notice is delivered at licenses/arxiv-2412.11697-CC-BY-4.0.txt. Characters: every retained excerpt is pure ASCII, so the delivered engine, XeLaTeX with its default Latin Modern text font, reports no missing character.
\begin{equation}\label{eq:PiFIO}
\Pi(x,y)=\int_0^{+\infty} 
e^{it\varphi(x,y)}s(x,y,t)dt+F(x,y),
\end{equation}

\begin{equation}\label{e-gue241220yydp}
\varphi(x,y)=x_{2n+1}-y_{2n+1}+\frac{i}{2}\sum_{j=1}^n\left[|z_j-w_j|^2+(\overline{z}_jw_j-z_j\overline{w}_j)\right]+O\left(|(x,y)|^4\right), 
\end{equation}

\begin{lemma}
With the notations used above, for $s_0(x,y)$ in \eqref{eq:PiFIO}, we
can take $s_0(x,y)$ so that $s_0(x,y)$ is independent of $y_{2n+1}$, 

\begin{equation}\label{e-gue241128yyd}
s_0(x,y)=s_0(x,y'),\ \ \mbox{for all $(x,y)\in D\times D$},
\end{equation}
where $y'=(y_1,\ldots,y_{2n})$,
\begin{equation}\label{e-gue241120yydg}
s_0(0,0)=\frac{1}{2\pi^{n+1}}\frac{dV_\xi}{dV}(0)=\frac{1}{2\pi^{n+1}}(e^{-2(n+1)f})(0),
\end{equation}

\begin{equation}\label{e-gue241120yydu}
\begin{split}
&\frac{\pr s_0}{\pr\ol z_j}(0,0)=\frac{\pr s_0}{\pr w_j}(0,0)=0,\ \ j=1,\ldots,n,\\
&\frac{\pr s_0}{\pr z_j}(0,0)=\frac{1}{2\pi^{n+1}}\frac{\pr}{\pr z_j}(e^{-2(n+1)f})(0),\ \ j=1,\ldots,n,\\
&\frac{\pr s_0}{\pr\ol w_j}(0,0)=\frac{1}{2\pi^{n+1}}\frac{\pr}{\pr\ol w_j}(e^{-2(n+1)f})(0),\ \ j=1,\ldots,n,\\
&\frac{\pr^2s_0}{\pr\ol z_j\pr z_\ell}(0,0)=\frac{\pr^2s_0}{\pr\ol w_j\pr w_\ell}(0,0)=0,\ \ j, \ell=1,\ldots,n,
\end{split}
\end{equation}
and 
\begin{equation}\label{e-gue241125ycd}
(\mathcal{T}_xs_0)(0,0)=(\mathcal{T}_ys_0)(0,0)=0,
\end{equation}
where $\frac{\pr}{\pr w_j}=\frac{1}{2}(\frac{\pr}{\pr y_{2j-1}}-i\frac{\pr}{\pr y_{2j}})$, $\frac{\pr}{\pr\ol w_j}=\frac{1}{2}(\frac{\pr}{\pr y_{2j-1}}+i\frac{\pr}{\pr y_{2j}})$, $j=1,\ldots,n$, $f(x):=\frac{1}{2(n+1)}\log \frac{dV(x)}{dV_\xi(x)}\in\cC^\infty(X)$. 
\end{lemma}

\begin{proof}
By using integration by parts in $t$, we can take $s_0$ so that $s_0$ is independent of $y_{2n+1}$.

We assume that $s_0$ is independent of $y_{2n+1}$. Since $\varphi$ is independent of $y_{2n+1}$, $\ddbar_{b,x}\varphi(x,y)$ vanishes to infinite order at $x=y$. From this observation and $s_0$ is independent of $y_{2n+1}$, we conclude that 

\begin{equation}\label{e-gue241207yydp}
\mbox{$\ddbar_{b,x}s_0(x,y')$ vanishes to infinite order at $x=y$}.
\end{equation}
From \eqref{e-gue241207yydp}, we get 
\begin{equation}\label{e-gue241207yydq}
\frac{\pr s_0}{\pr\ol z_j}(0,0)=0,\ \ j=1,\ldots,n.
\end{equation} 

From $\ddbar_{b,x}\int e^{-it\ol\varphi(y,x)}\ol s(y,x,t)dt\equiv0$, we get 
\begin{equation}\label{e-gue241207yydr}
\mbox{$\ddbar_{b,x}(-\ol\varphi(y,x))+g(x,y)\ol\varphi(y,x)$ vanishes to infinite order at $(0,0)$},
\end{equation}
where $s(x,y,t)$ is as in \eqref{eq:PiFIO}  and $g$ is a smooth function defined near $(0,0)$. From \eqref{e-gue241220yydp} and \eqref{e-gue241207yydq}, we can check that 
\begin{equation}\label{e-gue241207yyds}
g(x,y)=O(\abs{(x,y)}^2).
\end{equation}

By using integration by parts in $t$, we get 

\begin{equation}\label{e-gue241207ycdk}
\begin{split}
&\ddbar_{b,x}\int e^{-it\ol\varphi(y,x)}\ol s(y,x,t)dt\\
&=\int e^{-it\ol\varphi(y,x)}(-(n+1)g(x,y)\ol s_0(y,x)+\ddbar_{b,x}\ol s_0(y,x))t^n+O(t^{n-1})dt.
\end{split}
\end{equation}
From \eqref{e-gue241207ycdk}, we see that 

\begin{equation}\label{e-gue241207ycdo}
\mbox{$(-(n+1)g(x,y)\ol s_0(y,x)+\ddbar_{b,x}\ol s_0(y,x))-h(x,y)\ol\varphi(y,x)$ vanishes to infinite order at $(0,0)$}, 
\end{equation}
where $h$ is a smooth function defined near $(0,0)$. From \eqref{e-gue241207yyds} and \eqref{e-gue241207ycdo}, it is not difficult to see that 
\begin{equation}\label{e-gue241207ycdg}
\frac{\pr s_0}{\pr w_j}(0,0)=\frac{\pr^2 s_0}{\pr w_j\pr\ol w_j}=0,\ \ j=1,\ldots,n.
\end{equation}

From \eqref{e-gue241207yydq} and \eqref{e-gue241207ycdg}, we get the first equation in \eqref{e-gue241120yydu}. 

From \eqref{e-gue241207yydq}, \eqref{e-gue241207ycdg} and notice that $s_0(x,x')=\frac{1}{2\pi^{n+1}}(e^{-2(n+1)f})(x)$, we get the second and third equation in \eqref{e-gue241120yydu}.

From \eqref{e-gue241207yydp} and \eqref{e-gue241207ycdg}, we get the last equation in  \eqref{e-gue241120yydu}.

Finally, from $s_0(x,x')=\frac{1}{2\pi^{n+1}}(e^{-2(n+1)f})(x)$ and $f$ is $\mathcal{T}$-invariant, we get \eqref{e-gue241125ycd}.
\end{proof}

\end{document}
